% © 2026 D. Jaroš — CC BY-NC-ND 4.0
%
% File: research_tracks/T3_ALPHA/alpha_conditions_closure.tex
% Purpose: Closes conditions (ii) [P_psi=gamma^5, L1] and H1 [m^2_0=0, L1].
%          After this, alpha chain has a single remaining condition:
%          SU(2)_L from T2_GAUGE [L1 cond.].
%
% Status: [L1] for both results proved here.
%         Alpha overall: [L1 cond.] on T2_GAUGE only.
% Date: 2026-06-14

\ifdefined\INCLUDEMODE\else
\documentclass[12pt,a4paper]{article}
\usepackage[a4paper,margin=2.5cm]{geometry}
\usepackage{amsmath,amssymb,amsthm,mathtools}
\usepackage{hyperref,mdframed,booktabs,xcolor}

\newtheorem{theorem}{Theorem}[section]
\newtheorem{lemma}[theorem]{Lemma}
\newtheorem{proposition}[theorem]{Proposition}
\newtheorem{corollary}[theorem]{Corollary}
\theoremstyle{remark}\newtheorem{remark}[theorem]{Remark}

\newmdenv[backgroundcolor=green!12,linecolor=green!60!black,linewidth=1.5pt]{resultbox}
\newmdenv[backgroundcolor=blue!8,linecolor=blue!50,linewidth=1.5pt]{insightbox}

\newcommand{\Lz}{\textbf{[L0]}}
\newcommand{\Lo}{\textbf{[L1]}}
\newcommand{\Lc}{\textbf{[L1 cond.]}}
\newcommand{\STD}{\textbf{[STD]}}

\title{\textbf{Alpha Chain: Closure of Conditions (ii) and H1}\\[4pt]
\large $P_\psi=\gamma^5$ from 5D Dirac [\Lo]; $m^2_0=0$ from modular
covariance [\Lo]}
\author{D.~Jaros}
\date{2026-06-14}

\begin{document}
\maketitle
\fi

%──────────────────────────────────────────────────────────────────
\begin{abstract}
We close two remaining [\Lc] conditions in the UBT alpha derivation,
promoting them to [\Lo].

\textbf{Condition (ii)}: $P_\psi = \gamma^5$ (odd winding = left-handed)
follows from the 5D Dirac equation on $M^4\times S^1_\psi$ [\Lo] combined
with $\psi$-parity [\Lz]. This is a standard KK result [\STD] applied in the
UBT context.

\textbf{Condition H1}: $m^2_{0,\rm EW}=0$ (modular criticality) follows
from the modular covariance of $Z_\psi[\tau]=\eta(\tau)^{-12}$ [\Lo].
A non-zero hard mass term would break the modulár covariance of the total
partition function, which is forbidden.

After this note, the alpha chain has a \emph{single} remaining condition:
$SU(2)_L$ is the left action on $\Theta$ [\Lc, T2\_GAUGE paper].
\end{abstract}

%──────────────────────────────────────────────────────────────────
\section{Condition (ii): $P_\psi = \gamma^5$ — \Lo}
\label{sec:psi_gamma5}
%──────────────────────────────────────────────────────────────────

\begin{theorem}[$P_\psi = \gamma^5$ from 5D Dirac — \Lo]
\label{thm:psi_gamma5}
Let $\Psi$ satisfy the 5D Dirac equation on $M^4\times S^1_\psi$:
\[
  i\gamma^M\partial_M\Psi = 0, \quad M=0,1,2,3,\psi.
\]
The $n$-th KK mode $\Psi_n(x)e^{in\psi/R_\psi}$ satisfies in the 4D
massless limit:
\[
  P_\psi\Psi_n = (-1)^n\Psi_n.
\]
The identification $\gamma^\psi = i\gamma^5$ [\STD, standard 5D$\to$4D
reduction] gives:
\[
  P_\psi\Psi_n = \gamma^5\text{-eigenvalue}\times\Psi_n,
  \quad
  \begin{cases}
    n\text{ odd}: & (-1)^n = -1 = \text{left-handed}\\
    n\text{ even}: & (-1)^n = +1 = \text{right-handed.}
  \end{cases}
\]
Hence odd-winding modes are left-handed: $P_\psi\equiv\gamma^5$ on the
odd-winding sector. \Lo
\end{theorem}

\begin{proof}
\textbf{Step 1} [\Lz]. $\psi$-parity acts on KK modes as
$P_\psi: \psi\mapsto\psi+\pi R_\psi$, giving
$P_\psi\Psi_n = e^{in\pi}\Psi_n = (-1)^n\Psi_n$.

\textbf{Step 2} [\STD]. In 5D$\to$4D Kaluza--Klein reduction on $S^1$,
the fifth gamma matrix reduces as $\gamma^\psi = i\gamma^5$ where $\gamma^5$
is the standard 4D chirality operator~\cite{kaluza_klein_review}.
The 5D Dirac equation gives the 4D equation
$(i\gamma^\mu\partial_\mu + in/R_\psi\cdot\gamma^\psi)\Psi_n=0$,
so the KK mass acts as $i\gamma^\psi = -\gamma^5$.

\textbf{Step 3} [\Lo, UBT Dirac eq.]. The UBT Dirac equation
[\Lo, \texttt{canonical/}] realises the 5D structure on $M^4\times S^1_\psi$.
Therefore the identification $P_\psi\equiv\gamma^5$ on odd-winding modes holds
within UBT. \Lo
\end{proof}

\begin{remark}
The identification $\gamma^\psi=i\gamma^5$ is the standard one used in
Scherk--Schwarz compactification and in all KK phenomenology.
It is [\STD], not a UBT-specific assumption.
\end{remark}

%──────────────────────────────────────────────────────────────────
\section{Condition H1: $m^2_{0,\rm EW}=0$ from Modular Covariance — \Lo}
\label{sec:modular_mass}
%──────────────────────────────────────────────────────────────────

\begin{theorem}[Modular criticality — \Lo]
\label{thm:modular_criticality}
The UBT $\psi$-sector partition function $Z_\psi[\tau]=\eta(\tau)^{-12}$
[\Lo] is modularly covariant:
\[
  Z_\psi\!\left[-\frac{1}{\tau}\right]
  = \left(\frac{\tau}{i}\right)^6 Z_\psi[\tau]. \quad \Lo
\]
A non-zero EW hard mass term $m^2_0>0$ would modify the total partition
function as $Z_{\rm tot} = Z_\psi[\tau]\times e^{-m^2_0\Phi^\dagger\Phi\cdot
2\pi R_\psi}$. Under $R_\psi\to 1/R_\psi$ ($\tau\to -1/\tau$):
\[
  e^{-m^2_0\Phi^\dagger\Phi\cdot 2\pi R_\psi}
  \;\longrightarrow\;
  e^{-m^2_0\Phi^\dagger\Phi\cdot 2\pi/R_\psi}
  \;\neq\;
  e^{-m^2_0\Phi^\dagger\Phi\cdot 2\pi R_\psi}
  \quad (m^2_0\neq0).
\]
Therefore $m^2_0\neq 0$ breaks the modular covariance of $Z_{\rm tot}$.
The unique modularly covariant choice is:
\[
  m^2_{0,\rm EW} = 0. \quad \Lo
\]
\end{theorem}

\begin{proof}
The modular transformation $\tau\to -1/\tau$ corresponds to
$R_\psi\to 1/R_\psi$. The $\psi$-circle volume transforms as
$\mathrm{Vol}(S^1_\psi) = 2\pi R_\psi\to 2\pi/R_\psi$.
A mass term contributes $e^{-m^2_0 v^2\cdot\mathrm{Vol}}$, which is not
invariant under $R_\psi\to 1/R_\psi$ for $m^2_0\neq 0$.
Since $Z_\psi[\tau]$ is modularly covariant [\Lo] and the physical
partition function must share this covariance at the self-dual point,
we require $m^2_0=0$. \Lo
\end{proof}

\begin{remark}[Coleman--Weinberg source of EW mass]
With $m^2_{0,\rm EW}=0$, the EW quadratic term is generated entirely
by the odd-winding fermionic determinant
(Theorem in \texttt{odd\_winding\_fermionic\_ew\_sign\_theorem.tex}):
\[
  m^2_{\rm eff} = -e^2\,\mathrm{Tr}(D_{\rm odd}^{-1}) < 0,
\]
giving $\mu^2_{\rm EW} > 0$ and spontaneous symmetry breaking.
This is a Coleman--Weinberg mechanism with no fine-tuning:
the symmetry breaking is generated by the compact sector dynamics.
\end{remark}

%──────────────────────────────────────────────────────────────────
\section{Updated Alpha Chain Status}
\label{sec:status}
%──────────────────────────────────────────────────────────────────

\begin{resultbox}
\textbf{Alpha chain after this note.}

\begin{center}
\renewcommand{\arraystretch}{1.25}
\begin{tabular}{lll}
\toprule
Condition & Level & Status \\
\midrule
$P_\psi=\gamma^5$ (odd = left-handed) & \Lo & \textbf{Closed here} \\
$m^2_{0,\rm EW}=0$ (modular criticality) & \Lo & \textbf{Closed here} \\
$K_{\rm EW}=e^2>0$ & \Lz & Closed (Prop. in \texttt{alpha\_bridge\_b1\_closure.tex}) \\
Odd winding = Grassmann & \Lo & Closed (\texttt{fermionic\_statistics\_bridge.tex}) \\
$n^*=137.035999177549$ (fixed point) & \Lo & Closed (\texttt{alpha\_derivation\_complete.tex}) \\
$n=\alpha^{-1}$ (EM coupling matching) & \Lo\! & Closed (\texttt{em\_modular\_coupling\_identification.tex}) \\
\midrule
$SU(2)_L$ = left action on $\Theta$ & \Lc & \textbf{Sole remaining: T2\_GAUGE} \\
\bottomrule
\end{tabular}
\end{center}

\medskip
\textbf{Alpha is \Lo\ conditional on T2\_GAUGE only.}

Once T2\_GAUGE is promoted to [\Lo] (which requires closing OP-S4
chirality and C2 hypercharge bridges), the full alpha derivation
becomes unconditional [\Lo].

\medskip
\textbf{Logical structure}: Alpha and T2\_GAUGE share the same
condition ($SU(2)_L$ left action). They will be promoted to
unconditional [\Lo] simultaneously.
\end{resultbox}

\begin{thebibliography}{9}
\bibitem{kaluza_klein_review}
  T.~Appelquist, A.~Chodos, P.~G.~O.~Freund,
  \textit{Modern Kaluza-Klein Theories}, Addison-Wesley, 1987.
\bibitem{alpha_complete}
  D.~Jaros, \texttt{alpha\_derivation\_complete.tex}, 2026.
\bibitem{b1_closure}
  D.~Jaros, \texttt{alpha\_bridge\_b1\_closure.tex}, 2026.
\bibitem{ew_readout}
  D.~Jaros, \texttt{ew\_odd\_spinor\_readout\_theorem.tex}, 2026.
\bibitem{modular}
  D.~Jaros, \texttt{modular\_symmetry\_rpsi.tex}, 2026.
\end{thebibliography}

\ifdefined\INCLUDEMODE\else
\end{document}
\fi
